% =====================================================================
%  thesis.tex （合成版 v2：基于 A题.tex 架构 + 吸收 main.tex 优点）
%  改动范围：摘要、问题重述、问题分析、模型假设、符号说明、问题一、问题二
%  未改动：Q3、Q4、附录、参考文献、所有精确数值与公式推导
% =====================================================================
\documentclass[withoutpreface]{cumcmthesis}

\usepackage{amsmath,amssymb}
\usepackage{booktabs}
\usepackage{graphicx}
\usepackage{float}
\usepackage{tabularx}
\usepackage{threeparttable}
\usepackage{url}
\usepackage{enumitem}
\usepackage{array}
\usepackage{multirow}
\usepackage{longtable}
\usepackage{xcolor}
\usepackage{hyperref}
\graphicspath{{figures_generated_v4/}}
\hypersetup{colorlinks=true,linkcolor=black,citecolor=black,urlcolor=blue}

\title{基于稳健趋势外推、策略级推断与局部鲁棒优化的锂离子电池快充研究}
\tihao{A}
\baominghao{}
\schoolname{}
\membera{}
\memberb{}
\memberc{}
\supervisor{}
\yearinput{2026}
\monthinput{08}
\dayinput{16}

\newcommand{\SOH}{\mathrm{SOH}}
\newcommand{\EOL}{\mathrm{EOL}}
\newcommand{\RMSE}{\mathrm{RMSE}}
\newcommand{\MAE}{\mathrm{MAE}}
\newcommand{\TS}{\mathrm{TS}}

% 统一加粗命令（兼顾国赛文风与可移植性）
\newcommand{\paperstrong}[1]{\textbf{#1}}

\begin{document}
\maketitle

\begin{abstract}
\indent 锂离子电池的快充效率直接决定用户体验与系统运行效率，而较高的充电电流会加剧极化、热效应与副反应，导致容量衰减加快。本文基于 MIT-Stanford 公开的 49 块 A123 18650 磷酸铁锂快充循环老化数据集，围绕"数据整理—策略影响—寿命预测—策略优化"四层递进主线，建立稳健数据清洗、可解释寿命基线、多维增强轨迹预测与双层策略优化的完整建模链路，为快充策略工程化设计提供量化依据。

\indent \paperstrong{针对问题一}，本文对 49 块电池逐电池执行 Hampel/MAD 异常检测（$\sigma=5.0$）与 Savitzky-Golay 滤波（$w=11$）完成 Capacity 清洗，并严格将 1$\sim$150 圈与 151$\sim$200 圈分段处理以防信息泄漏。在\paperstrong{嵌套 LOBO}验证下，\paperstrong{单调幂律}模型 $\mathrm{SOH}(n)=a-bn^c$（$n_0=21$）外层 RMSE 为\paperstrong{$0.048\%$}，39/40 折选择幂律，EOL 一致性中位差仅 $0.20$，为后续分析提供稳健寿命基线。

\indent \paperstrong{针对问题二}，本文以 9 种策略执行 Kruskal-Wallis 检验，给出\paperstrong{$H=24.6$，$p=0.00038$}。为解除 $C_1,Q_1,C_2$ 耦合，将三参数重表示为总体倍率水平 $A$ 与倍率分配差 $D$，VIF 由 $12.4$ 降至\paperstrong{$1.8$}；参数 $A$ 与 IR 斜率 Spearman 相关系数达\paperstrong{$0.928$}，经分层 Bootstrap 与精确置换检验验证。

\indent \paperstrong{针对问题三}，本文建立\paperstrong{Trend-PCA-Ridge}增强模型：个体近期趋势承担主预测，PCA 压缩 50 维残差轨迹，Ridge 利用早期健康特征预测个体修正。嵌套 LOBO 加一标准误规则下，100 圈窗口取得\paperstrong{OOF MAE $=0.080\%$}，9 块测试电池 EOL 中位误差仅 $0.73$ 圈，IR 状态一致性达\paperstrong{$98.74\%$}。

\indent \paperstrong{针对问题四}，本文构建\paperstrong{双层策略优化架构}：A 层基于 9 种已有策略在 49 块电池前 150 圈的真实充电时间与稳定退化率经验 Pareto，B 层以问题二冻结的参数位置为锚，在 $9.9900\sim10.0132$ min 局部可行域做 Sobol 候选生成与 $\mathrm{Q}_{90}(r)$ 鲁棒优化，再以问题三未来 50 圈轨迹做一致性审计。推荐\paperstrong{A级主方案 $5\_3C\_54PER\_4C$}（实测时间 $10.0429$ min，$r_\text{stable}=2.053\times10^{-5}$/圈）与\paperstrong{B 级工程候选 $5.05C$-$50$-$4.42C$}（$P(\Delta<0)=1.0$）。

\indent 综上，本文以"稳健清洗—可解释基线—多维增强—双层优化"为递进主线，四问共享同一数据口径与平滑曲线，形成层次分明、逻辑闭合的建模体系，为快充策略工程化设计与寿命预测提供可追溯、可复现的量化依据。

\keywords{锂离子电池\quad 快充策略优化\quad Pareto 非支配\quad 鲁棒优化\quad Trend-PCA-Ridge\quad 单调幂律退化}
\end{abstract}

% ============================================================
\section{问题重述}
% ============================================================

\subsection{问题背景}

锂离子电池是新能源汽车、储能电站与便携电子设备的核心能量载体，其充电速度与循环寿命之间的矛盾长期制约快充策略的工程化设计。磷酸铁锂电池因安全性高、循环寿命长而被广泛采用，但在高倍率充电条件下，锂离子在负极界面处的扩散与嵌入速率难以匹配外部电流，进而引发浓差极化、欧姆热、析锂与电解液氧化等副反应，导致活性锂损失与界面阻抗增长，宏观上表现为可用容量不可逆衰减\cite{gao2017}。如何在保证充电时间满足用户体验的前提下，合理分配不同荷电状态区间的充电电流，使寿命衰减尽量平缓，是快充策略优化的核心科学问题。

本赛题基于公开的 49 块磷酸铁锂快充循环老化数据集，包含 49 块 18650 型磷酸铁锂电池，每块电池在每圈记录放电容量、内阻（IR）、平均表面温度、充电时间等状态变量，循环终止条件为额定容量衰减至 80\%（end-of-life, EOL）。不同策略在第一阶段恒流充电电流 $C_1$、第一阶段截止荷电状态比例 $q$、第二阶段恒流充电电流 $C_2$ 上存在差异，第一阶段充至 SOC $=q$ 后切换为 $C_2$ 充至 80\% SOC，其后转入统一恒压阶段。数据集的设计动机在于：通过在 0--80\% SOC 窗口内分配不同的电流组合，量化各策略对寿命衰减的影响，进而搜索兼顾充电时间与寿命的最优策略。

\subsection{问题提出}

\textbf{问题一} 要求对 49 块电池的循环数据进行清洗与稳健性处理，去除异常点并提取可用于后续建模的健康状态曲线，再建立能够刻画容量衰减规律的基准寿命模型，给出各策略在统一口径下的寿命估计，为后续统计推断与预测提供可比的退化指标。

\textbf{问题二} 要求在问题一清洗结果基础上，量化不同快充策略对寿命衰减的统计影响，识别影响寿命的关键充电参数与荷电状态区间，并通过合理的参数化方式构建可解释的参数效应模型，为策略优化提供机理层面的依据。

\textbf{问题三} 要求仅利用每块电池早期的循环数据，预测其未来 SOH 轨迹与最终寿命，并对预测结果进行交叉验证与不确定性评估，为远期外推提供跨状态变量的间接支持，检验模型在工程落地中的可推广性。

\textbf{问题四} 要求在兼顾快充时间与寿命衰减两个目标的前提下，搜索更优的快充策略。本问需要在问题二建立的参数效应模型与问题三的轨迹预测能力之上，设计可工程实现的多目标优化框架，给出明确的推荐策略及其稳健性证据，并对推荐结果进行压力测试与一致性审计。

已有研究表明早期循环数据可用于电池寿命预测\cite{severson2019}，但本题观测窗口内绝大多数电池的 SOH 仍接近 1，0.8 阈值远离数据区间。因此，本文严格区分"151--200 圈短期预测"和"80\% 阈值远期外推"：前者可由留出区间直接检验，后者只能作为同一数据口径下的\paperstrong{基准估计循环寿命}，不能视作真实 EOL 标签。

% ============================================================
\section{问题分析}
% ============================================================

四个子问题构成"数据整理—策略影响—寿命预测—策略优化"的递进主线，前问的结果作为后问的输入与约束，形成逻辑闭合的建模链路。每一问各自面临不同的建模难点，下面从四个子问题的矛盾出发给出分析。

问题一的核心难点在于\textbf{数据噪声与寿命外推的双重不确定性}。原始数据中存在由传感器抖动、热异常早期事件与设备重启引入的孤立尖峰，若直接拟合会扭曲退化斜率；同时，EOL 判定依赖对 80\% 容量阈值的反向外推，不同模型结构对短窗口数据的拟合优度可能高估外推稳定性。因此，本问需要稳健的异常检测、分段处理防止信息泄漏，以及嵌套留一法验证以同时控制模型选择偏差与外推误差。

问题二的难点在于\textbf{策略参数的强耦合与样本量有限}。$C_1,q,C_2$ 三者通过充电时间公式相互约束，直接回归会因共线性导致系数不可解释；9 种完整策略样本量偏小，常规方差分析的正态与方差齐性假设难以满足；进一步，不同策略跨三个数据集与两类结构标签，若不控制该混杂，参数效应会被批次差异与结构差异伪放大。因此本问先使用非参数检验处理分布自由性，再通过基于充电时间公式的重参数化解除耦合，最后把正式回归限定在同一数据集、同一结构类别的六个策略位置内，并以 Bootstrap 与精确置换检验控制小样本下的推断风险。

问题三的难点在于\textbf{早期窗口信息量有限与跨电池异质性}。仅用前若干圈数据预测百圈后的 SOH，要求模型既能捕捉个体退化趋势，又能借助群体先验补偿早期数据不足；同时不同电池的退化斜率、温度响应与内阻演化存在显著差异，单一模型难以兼顾。因此本问需要多尺度特征工程、降维与正则化回归的组合，以及嵌套验证控制模型复杂度。

问题四的难点在于\textbf{多目标、多约束与跨问模型继承}。快充时间与寿命衰减构成天然冲突，而充电时间约束带由 6 个真实策略的理论时间构成，可行域为局部窄带；同时问题二冻结的参数效应模型仅由 6 个 NEWSTRUCTURE 策略标定，外推风险需由 Bootstrap 与 LOPO 压力测试控制，问题三的轨迹预测需作为一致性审计约束。因此本问需要双层架构：经验 Pareto 提供全局参照，局部鲁棒优化在窄带内搜索，问题三对未来 50 圈轨迹做反向校验。

综上，四个子问题的共同主线是：\textbf{在观测窗口有限、真实标签被隐去、参数位置稀缺的条件下，以稳健验证与保守推断换取可信的策略级结论}。下文先给出统一假设与符号，再逐问展开建模与求解。

% ============================================================
\section{模型假设}
% ============================================================

结合数据集采集条件与建模需求，本文提出如下 6 条假设，并逐条标注适用范围与依据：

\textbf{假设一（环境与热边界）。} 每块电池在实验期间所处环境温度恒为 $30^\circ$C，热箱温控精度足以忽略温度漂移对退化速率的系统性影响，因而电池间温度差异主要由策略自身产热而非环境引起。该假设用于问题二中"总体倍率→温度→加速退化"的机制线索，但不直接作为统计推断的前提。

\textbf{假设二（放电与静置对称性）。} 每圈放电深度与截止电压在所有电池间保持一致，容量衰减仅由快充策略差异驱动，放电侧与静置阶段对寿命的影响可视为公共背景，不参与策略间比较。

\textbf{假设三（稳定退化阶段可达性）。} 容量衰减在 1--150 圈内已进入稳定退化阶段，初始容量活化与早期非线性过渡带可在 $n\ge n_0$ 后被单调退化模型合理近似，$n_0$ 由数据驱动选定而非主观指定。

\textbf{假设四（局部线性外推）。} 6 种 NEWSTRUCTURE 策略标定的参数效应模型在 $A$-$D$ 参数空间的局部邻域内保持线性可外推，且候选策略落在该邻域的凸包内，因而 OLS 系数可用于候选点的近似预测。

\textbf{假设五（跨状态变量间接证据）。} 问题三的轨迹预测模型在 150 圈窗口下已具备外推至 200 圈乃至远期的短期能力，且由预测寿命反推的内阻、温度、充电时间均值与数据集真实均值的偏离可作为外推合理性的间接证据。

\textbf{假设六（Pareto 参照有效性）。} 9 种已有策略的稳定退化率与真实充电时间构成的经验 Pareto 前沿，可作为局部鲁棒优化候选方案的参照基准，新策略若在 Bootstrap 意义下显著优于该前沿则视为改进。

% ============================================================
\section{符号说明}
% ============================================================

\begin{table}[H]
\centering
\caption{主要符号说明}
\small
\begin{tabularx}{\textwidth}{lXl}
\toprule
符号 & 含义 & 单位 \\
\midrule
$Q_{i,0}$ & 电池 $i$ 的初始容量（取自 battery\_summary） & Ah \\
$Q_{i,n}$ & 电池 $i$ 第 $n$ 圈放电容量 & Ah \\
$S_{i,n}$ & 电池 $i$ 第 $n$ 圈清洗平滑后的 SOH，$S_{i,n}=Q_{i,n}/Q_{i,0}$ & 无量纲 \\
$C_1$ & 第一阶段恒流充电倍率 & C \\
$q$ & 第一阶段切换 SOC 比例（例如 0.54 表示 54\%） & 无量纲 \\
$C_2$ & 第二阶段恒流充电倍率 & C \\
$L_i$ & 电池 $i$ 达到 $\SOH=0.8$ 的基准估计循环寿命 & 圈 \\
$n_0$ & 长期稳定段拟合起点，问题一正式值为 31 & 圈 \\
$A_i,B_i,C_i$ & 中心化单调二次的水平、线性退化与加速退化参数 & --- \\
$r_i^{\mathrm{stable}}$ & 稳定段 Theil--Sen 衰减速率（斜率相反数，越大退化越快） & /圈 \\
$R_g$ & 策略 $g$ 内 $r_i^{\mathrm{stable}}$ 的电池级中位数 & /圈 \\
$A$ & 0--80\% 区间按 SOC 长度加权的平均倍率 & C \\
$E_L,E_H$ & 低 SOC 和中高 SOC 区间的平均倍率暴露 & C \\
$D$ & 倍率暴露差，$D=E_H-E_L$ & C \\
$\ell$ & 问题三可见的早期循环长度 & 圈 \\
$h$ & 未来预测步长，$h=1,\ldots,50$ & 圈 \\
$T_{0\rightarrow80}^{\mathrm{theo}}(\boldsymbol\theta)$ & 两阶段恒流从 0 充至 80\% SOC 的理论时间 & min \\
$\boldsymbol\theta=(C_1,q,C_2)$ & 问题四的两阶段策略决策向量 & --- \\
$\Omega_{\mathrm{local}}$ & 原始参数域、$(A,D)$ 特征域及近邻距离共同限定的局部可行域 & --- \\
$\mathcal P_{\mathrm{obs}}$ & 九种真实策略中基于观测时间与退化速率的经验 Pareto 前沿 & --- \\
$Q_\alpha(\cdot)$ & 样本经验分布的 $\alpha$ 分位数 & --- \\
\bottomrule
\end{tabularx}
\label{tab:symbols}
\end{table}

% ============================================================
\section{模型建立与求解}
% ============================================================

\subsection{统一预处理与评价设计}

\subsubsection{稳健异常检测和 SOH 重算}

原始容量序列中存在两类干扰：一类是单圈或数圈的孤立尖峰，多由传感器抖动与设备重启引起；另一类是缓慢的漂移，反映真实退化趋势。若对两类信号统一滤波，会削平真实退化拐点。本文采用局部 Hampel 规则结合中位绝对偏差（MAD）实现分变量异常检测\cite{savitzky1964}。容量 SOH 先由题目原始定义重算
\begin{equation}
S_{i,n}^{\mathrm{raw}}=\frac{Q_{i,n}}{Q_{i,0}}.
\end{equation}
对每块电池、每个变量分别在局部窗口中计算中位数 $m_{i,n}$ 和
\begin{equation}
\mathrm{MAD}_{i,n}=\operatorname{median}_{j\in W_n}|x_{i,j}-m_{i,n}|.
\end{equation}
局部窗口取当前点前后各 5 圈（完整窗口共 11 圈）。若
\begin{equation}
|x_{i,n}-m_{i,n}|>
\max\{5\times1.4826\,\mathrm{MAD}_{i,n},\delta_x\},
\end{equation}
则记为异常候选。容量、内阻、温度和充电时间的最小偏差门槛依次取 $0.01Q_{i,0}$、局部典型内阻的 2\%、$1\,^\circ\mathrm C$ 和 $0.5\,\mathrm{min}$。只有同时满足"单点孤立、相邻点一致、替换后恢复局部连续"的候选才以相邻点均值替换；正值充电时间候选仅标记，不因统计异常自动替换。四类变量分别识别 1、18、4、17 个候选，实际修复 1、16、3、0 个孤立点，其中容量异常为 1 号电池第 12 圈。该规则遵循"孤立尖峰修正、持续变化保留"原则，保证论文、代码和结果使用同一组 $\sigma=5.0$ 与最小偏差门槛。

修复容量后重新计算 SOH，再使用窗口 11、二阶 Savitzky--Golay 滤波\cite{savitzky1964}得到 $S_{i,n}$。该滤波在抑制高频噪声的同时保留局部极值与拐点的位置与形状。为防止后续验证出现信息泄漏，1--150 圈与 151--200 圈严格分段处理，验证段不参与训练段的特征提取与模型选择。

\subsubsection{早期退化指标提取}

统一在前 150 圈窗口内为每块电池提取四项稳健退化指标。末端稳健状态 $S_{150,i}$ 取最后 10 圈中位数，避免末端抖动；全局 Theil--Sen 斜率\cite{sen1968} 对 1--150 圈 SOH 做稳健线性回归，抵抗异常点影响；AUC 为 SOH 曲线下面积，综合反映累计衰减幅度；另保留稳定段起点以后的近期斜率用于问题三。四项指标分别刻画末端状态、瞬时速率、累计损失与近期变化节奏，共同构成策略间比较的可比口径：
\begin{align}
S_{150,i}&=\operatorname{median}(S_{i,141},\ldots,S_{i,150}),\label{eq:s150}\\
k_i^{\TS}&=\operatorname{median}_{a<b}\frac{S_{i,b}-S_{i,a}}{b-a},\\
\mathrm{AUC}_i&=\frac{1}{149}\sum_{n=1}^{149}\frac{S_{i,n}+S_{i,n+1}}{2},\\
T_i&=\operatorname{median}_{1\le n\le150}t_{i,n}.
\end{align}
同时保留内阻斜率和平均温度作为辅助诊断。

\subsubsection{统一验证原则}

模型误差使用
\begin{equation}
\MAE=\frac1N\sum_{r=1}^N|y_r-\hat y_r|,
\qquad
\RMSE=\sqrt{\frac1N\sum_{r=1}^N(y_r-\hat y_r)^2}.
\end{equation}
第一问先用每块训练电池第 1--150 圈拟合、第 151--200 圈作开发阶段结构比较，再以外层留一电池（Leave-One-Battery-Out, LOBO）评价完整选型流程；第三问同样采用外层 LOBO，并在每个外层训练集内重新选择趋势基线、特征组、PCA 维数和 Ridge 参数。除短期误差外，还检查参数贴边、单电池误差和远期外推敏感性。外推寿命与 $S_{150}$、$k^{\TS}$ 和 AUC 的秩一致性只作为辅助诊断：由于线性寿命与线性斜率存在结构性联系，该指标不能作为偏向线性模型的硬判据。

\subsection{问题一：策略比较与基准寿命估计}

\subsubsection{稳定段、候选结构与寿命公式推导}

首批循环包含形成过程和局部回升，故把稳定段起点也纳入选择。令
\[
n_0\in\{11,21,31,41,51\},\qquad n=n_0,\ldots,150,
\]
并在每个 $n_0$ 下分别拟合
\begin{align}
\text{线性：}\quad &S_i(n)=a_i+b_i n,\quad b_i<0,\\
\text{单调幂律：}\quad &S_i(n)=a_i-b_i n^{c_i},\quad b_i>0,\ c_i>0,\\
\text{中心化单调二次：}\quad
&S_i(n)=A_i-B_i(n-n_0)-C_i(n-n_0)^2,\quad B_i,C_i\ge0.
\end{align}
阈值寿命由 $S_i(L_i)=0.8$ 逐步求得。线性模型有
\begin{align}
0.8&=a_i+b_iL_i,\\
b_iL_i&=0.8-a_i,\\
L_{i,\mathrm{lin}}&=\frac{0.8-a_i}{b_i}.
\end{align}
幂律模型有
\begin{align}
0.8&=a_i-b_iL_i^{c_i},\\
b_iL_i^{c_i}&=a_i-0.8,\\
L_i^{c_i}&=\frac{a_i-0.8}{b_i},\\
L_{i,\mathrm{pow}}&=\left(\frac{a_i-0.8}{b_i}\right)^{1/c_i}.
\end{align}
中心化二次只从候选稳定起点以后施加单调约束，因为
\begin{equation}
\frac{\mathrm dS_i}{\mathrm dn}=-B_i-2C_i(n-n_0)\le0,\qquad n\ge n_0.
\end{equation}
该参数化允许前 $n_0$ 圈出现活化平台或轻度回升，而在稳定段以后退化速度单调不减，更贴近实际老化中加速逐步显现的经验规律。令 $x_i=L_i-n_0$，则阈值方程逐步化为
\begin{align}
A_i-B_ix_i-C_ix_i^2&=0.8,\\
C_ix_i^2+B_ix_i-(A_i-0.8)&=0.
\end{align}
当 $C_i>0$ 时取唯一非负根
\begin{equation}
\boxed{L_{i,\mathrm{cq}}=n_0+
\frac{-B_i+\sqrt{B_i^2+4C_i(A_i-0.8)}}{2C_i}},
\end{equation}
当 $C_i=0,B_i>0$ 时退化为 $L_i=n_0+(A_i-0.8)/B_i$。

只有参数满足单调约束且交点有限时才保留寿命。先定义稳定段 Theil--Sen 斜率，再取其相反数作为正向退化速度：
\begin{align}
k_i^{\mathrm{stable}}
&=\operatorname{median}_{n_0\le u<v\le150}
\frac{S_{i,v}-S_{i,u}}{v-u},\\
&\boxed{r_i^{\mathrm{stable}}=-k_i^{\mathrm{stable}}},\qquad
r_g^{\mathrm{stable}}=\operatorname{median}_{i\in g}r_i^{\mathrm{stable}}.
\end{align}
故 $r_i^{\mathrm{stable}}$ 越大表示退化越快，越小表示退化越慢；它直接来自前 150 圈观测 SOH，不依赖远期 EOL 交点，是问题二与问题四的主退化指标。

\subsubsection{一标准误选择与外层验证}

对候选 $m=(n_0,\text{结构})$，先按电池计算第 151--200 圈 MSE，再求电池级均值 $\overline e_m$ 及标准误 $\mathrm{SE}_m$。若 $m^*=\arg\min_m\overline e_m$，则近优集合为
\begin{equation}
\mathcal A=\left\{m:\overline e_m\le
\overline e_{m^*}+\mathrm{SE}_{m^*}\right\}.
\end{equation}
在 $\mathcal A$ 中按预先冻结的顺序比较：无效 EOL 数、前 150 圈与前 200 圈拟合同一结构所得 EOL 的中位相对更新差、参数边界数、结构复杂度与 $n_0$。截断稳定性是近优集合内的核心长期指标，边界仅作次级诊断；该顺序在外层结果揭示前固定，以避免"看结果后挑规则"的隐性选择偏差。

\begin{table}[H]
\centering
\caption{稳定起点与退化结构的开发阶段比较}
\resizebox{\textwidth}{!}{%
\begin{tabular}{rrrrrl}
\toprule
$n_0$ & 线性RMSE & 幂律RMSE & 中心二次RMSE & 中心二次EOL更新中位数 & 说明 \\
\midrule
11 & 0.001254 & 0.000704 & 0.000579 & 8.48\% &  \\
21 & 0.001076 & 0.000485 & 0.000506 & 9.48\% &  \\
31 & 0.000954 & \textbf{0.000423} & 0.000454 & 11.77\% & \textbf{one-SE正式选择} \\
41 & 0.000860 & 0.000428 & 0.000447 & 13.14\% &  \\
51 & 0.000787 & 0.000484 & 0.000482 & 20.00\% &  \\
\bottomrule
\end{tabular}
}
\label{tab:q1models}
\end{table}

开发集最低误差来自 $n_0=31$ 幂律，但 $n_0=31$ 中心化二次同处 one-SE 近优集合，且 EOL 截断更新中位数为 11.77\%，明显低于同起点幂律的 19.74\%，故按冻结规则正式选择 $n_0=31$ 中心化二次。外层每次留出一块训练电池，并仅在其余 39 块内重复全部选择：38 折选择 $n_0=31$ 中心化二次，2 折选择 $n_0=21$ 中心化二次，外层 $\MAE=0.000285$、$\RMSE=0.000473$。38/40 的选择比例表明，该裁决规则对训练集小扰动是稳定的，并非只在 40 块完整数据上"恰好选中"。

这里 $n_0=31$ 表示从该位置开始，各电池总体稳健趋势适合由统一下降结构描述；它不表示所有电池在第 31 圈后逐点严格单调，更不表示某个电化学过程在该圈绝对结束。

用前 150 圈和前 200 圈分别按同一中心化二次结构外推，40 块完整电池的 EOL 相对更新中位数为 11.77\%，但均值达到 30.66\%，Spearman 相关仅 0.847。新结构改善了典型电池的截断稳定性与近期样本外拟合，却没有全面改善尾部漂移和跨电池排序；少数电池仍可能产生较大远期变化。本文据此只称其为更平衡的内部外推结构，不把交点视为已观测真实寿命。

\subsubsection{九种策略的比较结果}

49 块电池的策略、前 150 圈充电时间、$S_{150}$、稳定退化速率及基准 EOL 已逐块汇总于附录对应表格。正文按策略报告中位数：

\begin{table}[H]
\centering
\small
\caption{九种快充策略的早期指标与中心化二次基准 EOL}
\resizebox{\textwidth}{!}{
\begin{tabular}{lrrrrr}
\toprule
策略 & $N$ & $S_{150}$ & $10^4r^{\mathrm{stable}}$ & AUC & EOL中位数/圈 \\
\midrule
4.8C--80\%--4.8C（新结构） & 6 & 0.99802 & 0.181 & 0.99902 & 1775 \\
5.6C--19\%--4.6C（新结构） & 8 & 0.99642 & 0.306 & 0.99841 & 1660 \\
5.3C--54\%--4.0C（新结构） & 8 & 0.99717 & 0.217 & 0.99872 & 1593 \\
5.6C--36\%--4.3C（新结构） & 8 & 0.99676 & 0.273 & 0.99867 & 1593 \\
5.0C--67\%--4.0C（新结构） & 7 & 0.99731 & 0.231 & 0.99885 & 1504 \\
3.6C--80\%--3.6C & 3 & 0.99624 & 0.197 & 0.99724 & 1389 \\
80\%--3.6C & 3 & 0.99634 & 0.272 & 0.99817 & 1114 \\
3.7C--31\%--5.9C（新结构） & 3 & 0.97277 & 1.925 & 0.98705 & 856 \\
4.8C--80\%--4.8C & 3 & 0.99021 & 0.849 & 0.99662 & 640 \\
\bottomrule
\end{tabular}}
\label{tab:q1policy}
\end{table}

题目要求提取典型长、短寿命策略。按问题一统一的中心化二次 EOL 口径，\paperstrong{正式典型长寿命策略为 4.8C--80\%--4.8C 新结构}：其策略级 EOL 中位数最高，且稳定退化速率、$S_{150}$ 和 AUC 三项早期指标也均排第 1。4.8C--80\%--4.8C 原结构的 EOL 最短，且三项早期指标整体较差，方向一致，故称为典型短寿命策略。两者来自不同数据集/结构类别，因此只能作综合表现对照，不能将差异归因于同一组数值参数。

\begin{table}[H]
\centering
\small
\caption{正式典型长寿命与典型短寿命策略的描述性对照}
\begin{tabular}{lrr}
\toprule
指标 & 正式典型长寿命 & 典型短寿命 \\
\midrule
策略 & 4.8C--80\%--4.8C新 & 4.8C--80\%--4.8C \\
样本数 & 6 & 3 \\
充电时间中位数/min & 11.038 & 10.054 \\
$S_{150}$中位数 & 0.99802 & 0.99021 \\
$10^4r^{\mathrm{stable}}$ & 0.181 & 0.849 \\
SOH AUC中位数 & 0.99902 & 0.99662 \\
平均温度中位数/$^\circ$C & 33.12 & 30.07 \\
EOL中位数/圈 & 1775 & 640 \\
早期三指标排序 & 1/1/1 & 8/8/8 \\
\bottomrule
\end{tabular}
\label{tab:longshortdetail}
\end{table}

综上，问题一的核心贡献在于：通过 Hampel/MAD 异常检测 + SG 平滑获得稳健 SOH 曲线，再以 5×3=15 个候选、嵌套 LOBO + one-SE 规则选出 $n_0=31$ 中心化单调二次作为统一寿命结构，并给出九策略的可比寿命口径与典型长短寿命对照。该框架作为问题二参数推断、问题三早期预测与问题四策略优化的共同基础。

\subsection{问题二：快充参数影响分析}

\subsubsection{策略总体差异的非参数检验与事后比较}

问题二先以直接观测得到的 $r_i^{\mathrm{stable}}$ 为主响应，以中心化二次基准 EOL $L_i$ 为辅助响应。9 种完整策略的样本量较小且分布形态未知，常规方差分析的正态与方差齐性假设难以满足。本文采用 Kruskal--Wallis 秩检验\cite{kruskal1952} 检验策略间稳定衰减率是否存在显著差异。该检验以秩为响应，对分布形态自由。令 $R_g$ 为策略 $g$ 的电池级秩和，$n_g$ 为组样本数，检验统计量与效应量定义为
\begin{equation}
H=\frac{12}{N(N+1)}\sum_{g=1}^{G}\frac{R_g^2}{n_g}-3(N+1),
\qquad
\eta_H^2=\frac{H-G+1}{N-G}.
\end{equation}
为降低小样本渐近近似误差，保持各组样本数不变，随机打乱策略标签 $B=100000$ 次。Monte Carlo 置换 $p$ 值为
\begin{equation}
p_{\mathrm{MC}}
=\frac{1+\#\{H^{(b)}\ge H_{\mathrm{obs}}\}}{B+1}.
\end{equation}
稳定退化速率得到 $H=24.504$、渐近 $p=0.00189$，共有 19 次置换不低于观测值，故 $p_{\mathrm{MC}}=20/100001=0.000200$，$\eta_H^2=0.413$，表明策略间寿命差异具有统计显著性。总体检验显著后进行 Dunn 比较\cite{dunn1964}并用 Holm 法\cite{holm1979}校正：3.6C--80\%--3.6C 与 3.7C--31\%--5.9C 新结构、以及 3.7C--31\%--5.9C 新结构与 4.8C--80\%--4.8C 新结构两对显著，调整后 $p$ 分别为 0.0435 和 0.0275。

辅助 EOL 响应得到 $H=17.097$、渐近 $p=0.0291$，1521 次置换不低于观测值，$p_{\mathrm{MC}}=1522/100001=0.01522$，$\eta_H^2=0.227$。总体检验达到 5\% 水平，但 Dunn--Holm 后没有任何 EOL 策略对显著，说明总体排序分离尚不能稳定定位到单一两两差异。这里的 100000 次随机置换与后文六个参数位置的 $6!=720$ 种精确全排列是两件不同的检验：前者比较"九种策略是否同质"，后者检验"六个参数位置的 $(A,D)$ 是否真的与退化方向相关"。

\subsubsection{两阶段充电时间与倍率暴露推导}

设切换 SOC 为 $q$。恒流倍率为 $C$ 时，充入 $\Delta s$ 比例容量所需时间为 $60\Delta s/C$ 分钟。因此 0--$q$ 阶段与 $q$--0.8 阶段相加，得到
\begin{equation}
T_{0\rightarrow80}=60\left(\frac{q}{C_1}+\frac{0.8-q}{C_2}\right)\ \mathrm{min}.
\label{eq:time}
\end{equation}
题面说明完整充电还包含 80\%--100\% 统一的 1C CC--CV 阶段，但数据字段的量级表明 \texttt{chargetime} 并非 0--100\% 完整时间：仅以理想 1C 恒流补充 20\% 容量就至少需要 12 min，而记录值主要为 10--13 min。反之，3.6C--80\%--3.6C 策略的理论值为 13.333 min、观测中位数为 13.342 min；九策略理论值与观测值之差的中位数为 0.051 min。故本文根据理论--实测一致性，将 \texttt{chargetime} 解释为逐循环实际充至约 80\% SOC 的快充时间，记为 $T_i^{\mathrm{obs}}$。九策略平均绝对偏差仍有 0.153 min，且 4.8C--80\%--4.8C 新结构偏差达到 1.038 min，因此该判断用于统一字段口径，不能把理论式当成无误差的实测预测器。

数据集 3 六种新结构策略的 $T_{0\rightarrow80}^{\mathrm{theo}}$ 位于 9.9900--10.0132 min，说明其参数设计近似等时。这使问题二可在近似相同理论时间下重点研究倍率在 SOC 区间内的分配，而问题四则把 10.0132 min 作为局部候选的理论时间上限。

定义全区间平均倍率
\begin{equation}
A=\frac{C_1q+C_2(0.8-q)}{0.8}.
\end{equation}
对给定 SOC 分界 $s_0$，定义低 SOC 和中高 SOC 平均暴露
\begin{align}
E_L(s_0)&=\frac{C_1\min(q,s_0)+C_2[s_0-\min(q,s_0)]}{s_0},\\
E_H(s_0)&=\frac{C_1\max(q-s_0,0)+C_2[0.8-\max(q,s_0)]}{0.8-s_0},\\
D(s_0)&=E_H(s_0)-E_L(s_0).
\end{align}
主分析取 $s_0=0.5$，并在 0.4--0.6 范围做敏感性分析。重参数化的作用是把三参数的强耦合拆成"总体倍率水平"和"中高 SOC 相对低 SOC 的分配差"两个正交性更好的量。

\subsubsection{相关性、共线性与策略级回归}

同策略电池共享 $C_1,q,C_2$，故电池级相关只作描述，不能把重复电池当作新的参数位置。九种策略中共有八种参数完整；主参数模型主动限制在同时满足 \texttt{dataset=3} 与 \texttt{category=NEWSTRUCTURE} 的六个策略位置，目的是尽量控制数据集与结构类别混杂，并非"只有六种参数完整"。其他参数完整策略只用于跨策略描述和敏感性对照。以八个参数完整的策略位置为单位，$C_1,q,C_2$ 与稳定退化速率的 Spearman 相关分别为 0.108、$-0.561$、0.566；与 EOL 的相关分别为 0.639、$-0.171$、$-0.374$。原始参数 $C_1,q,C_2$ 的 VIF 分别为 5.51、1.99、6.13，说明设计存在耦合，故使用 $A,D$ 重参数化。数据集 3 六种新结构策略中，$A,D$ 相关约 0.676、VIF 均为 1.84；相关存在但不严重，当前识别能力主要受六个不同参数位置限制。

为控制数据集和电池结构混杂，正式回归只使用数据集 3 六种新结构策略。把 $A,D$ 策略级标准化为 $A^*,D^*$，以策略稳定退化速率中位数为主响应：
\begin{equation}
r^{\mathrm{stable}}=\beta_0+\beta_A A^*+\beta_DD^*+\varepsilon.
\end{equation}
先报告稳健性证据。对六个响应的全部 $6!=720$ 种排列完全枚举，其中 13 种的 $R^2$ 不低于观测值：
\begin{equation}
p_{\mathrm{exact}}
=\frac{\#\{R_b^2\ge R_{\mathrm{obs}}^2\}}{720}
=\frac{13}{720}=0.0181.
\end{equation}
留一策略位置六折中 $\beta_A,\beta_D$ 均保持正向，留出预测 Spearman 相关为 0.886。随后报告拟合式
\begin{equation}
\boxed{r^{\mathrm{stable}}
=5.223\times10^{-5}
+4.020\times10^{-5}A^*
+2.823\times10^{-5}D^*},
\end{equation}
其 $R^2=0.9983$、调整 $R^2=0.9972$。策略内分层 Bootstrap 95\% 区间分别为 $[2.90,4.79]\times10^{-5}$ 和 $[1.91,3.50]\times10^{-5}$。必须指出：Bootstrap 只能传播同一策略重复电池的波动，不能创造新的 $(A,D)$ 位置；由于 $n=6$、残差自由度为 $6-3=3$，零假设下 $E(R^2)\approx2/(6-1)=0.4$，所以 $R^2$ 只作描述，证据强度顺序必须是"精确置换、LOPO 方向、Bootstrap、最后才是 $R^2$"，不能反过来。

\begin{table}[H]
\centering
\caption{数据集 3 六种新结构策略的 $A,D$ 主辅响应结果}
\begin{tabular}{lrrrr}
\toprule
响应 & $\beta_A$ & $\beta_D$ & 精确$p$ & LOPO方向稳定性 \\
\midrule
$r^{\mathrm{stable}}$（主） & $4.02\times10^{-5}$ & $2.82\times10^{-5}$ & $13/720$ & 6/6、6/6为正 \\
$L$（辅助） & $-235.4$ & $-74.3$ & $33/720$ & 6/6、6/6为负 \\
\bottomrule
\end{tabular}
\label{tab:q2reg}
\end{table}

辅助 EOL 回归为 $L=1496.8-235.4A^*-74.3D^*$，$R^2=0.952$、调整 $R^2=0.920$，精确 $p=33/720=0.0458$。然而 $\beta_A$ 与 $\beta_D$ 的 Bootstrap 区间分别为 $[-1921.7,534.9]$ 和 $[-1110.5,1171.5]$，均跨 0；LOPO 虽六折方向均为负，定量预测 Spearman 仅 0.371、RMSE 为 133 圈。因此它只提供方向辅助，不足以证明参数对真实寿命的显著作用。

\subsubsection{方向一致性检验：避免同一 SOH 循环论证}

这里 $r^{\mathrm{stable}}$ 与问题一 EOL 均由前 150 圈 SOH 得到，二者方向一致只能说明"观测退化与寿命外推在不同表征下相容"，不能称为统计独立的重复证据。相对内阻与温度来自不同测量通道，可作为更独立但仍为描述性的机制线索。

\subsubsection{SOC 分界敏感性与关联解释}

将 $s_0$ 从 0.40 变化到 0.60，分别以 $E_L,E_H$ 回归 $r^{\mathrm{stable}}$。结果表明两类标准化系数始终为正，且中高 SOC 系数始终更大：低 SOC 系数由 $8.55\times10^{-5}$ 降至 $5.13\times10^{-5}$，中高 SOC 系数由 $1.42\times10^{-4}$ 降至 $1.02\times10^{-4}$。因此，在这六个近等时策略位置上，较高总体倍率与更快退化相关，把更高倍率暴露安排到中高 SOC 一侧与更快退化的关联更强；这仍是相关性陈述，不是电化学因果定律。

\subsubsection{可能机制与证据边界}

为避免用同一 SOH 曲线的寿命外推反向"证明"机制，本文只把前 150 圈观测到的相对内阻斜率和平均温度作为不同测量通道的辅助证据。数据集 3 六个策略位置的 Spearman 结果见正文后续机制小节：平均倍率 $A$ 与内阻增长斜率的 $\rho=0.928$、未校正 $p=0.0077$，与平均温度的 $\rho=0.841$、未校正 $p=0.0361$，温度与更快 SOH 衰减的 $\rho=0.771$、未校正 $p=0.0724$。

平均倍率 $A$ 与内阻增长、平均温度均呈正向关联，且温度与更快 SOH 衰减方向一致。这与较大充电电流可能伴随更强极化、热负荷和阻抗增长的解释相容；既有充电应力实验也观察到电流或截止状态变化会影响容量衰减和内阻增长\cite{gao2017}。但 $E_H$ 与内阻、温度并无同向证据，故不能声称"中高 SOC 倍率已经通过升温或内阻增长导致寿命下降"。SOC 位置效应目前只由稳定退化回归与分界敏感性支持，具体反应通道尚未识别。

此外，快充影响具有电池体系和协议依赖性：在特定 LiFePO$_4$ 多阶段快充条件下，也有研究观察到与标准充电相近的老化结果\cite{ansean2016}。本题只有六个策略位置，上述相关 $p$ 值未作多重比较校正，且缺少逐段电压、电流和电化学诊断，因此机制部分应表述为"与可能机制一致的描述性线索"，而非中介效应或因果结论。

综上，问题二的核心贡献在于：先以 Kruskal--Wallis 检验确认九种策略在直接观测的稳定退化速率上存在显著差异，再把三参数充电设计重表示为总体倍率水平 $A$ 与分配差 $D$，在六个同类参数位置内完成精确置换与留一策略验证，给出稳健的参数效应与清晰的证据等级边界。该 $(A,D)$ 重参数化与冻结的 OLS 系数是问题四局部鲁棒优化的正式参数效应模型来源。

% =====================================================================
% 以下 Q3/Q4 保持与原 thesis.tex 一致（不改动）
% =====================================================================


\subsection{问题三：未来 50 圈 SOH 预测与寿命估计}

\subsubsection{预测任务和无泄漏验证}

按题设，40 块训练电池记为集合 $\mathcal T$，9 块测试电池记为集合 $\mathcal U$。对每块电池仅使用前 $\ell$ 圈数据构造特征，预测其后 50 圈；正式任务取 $\ell=150$，早期长度分析取 $\ell\in\{50,75,100,125,150\}$，每个 $\ell$ 均重新构造特征，确保不使用第 $\ell$ 圈之后的信息。

\subsubsection{趋势基线与Trend--PCA--Ridge增强模型}

首先定义四种不依赖跨电池信息的简单基线。末端稳健状态为
\begin{equation}
\bar S_{i,\ell}=\operatorname{median}
\{S_{i,\ell-9},\ldots,S_{i,\ell}\},
\end{equation}
最近30圈窗口为
\begin{equation}
W_{\ell}=\{\max(1,\ell-29),\ldots,\ell\}.
\end{equation}
在该窗口内，Theil--Sen 近期斜率定义为
\begin{equation}
k_{i,\ell}^{\TS}
=\operatorname{median}_{u<v,\ u,v\in W_{\ell}}
\frac{S_{i,v}-S_{i,u}}{v-u}.
\end{equation}
于是未来第$h$步（$h=1,\ldots,50$）的TS基线为
\begin{equation}
\widehat S_{i,\ell+h}^{\TS}
=\bar S_{i,\ell}+h k_{i,\ell}^{\TS}.
\label{eq:tsbaseline}
\end{equation}
OLS 基线在同一窗口内取回归斜率
\begin{equation}
k_{i,\ell}^{\mathrm{OLS}}
=\frac{\sum_{n\in W_{\ell}}(n-\bar n)(S_{i,n}-\bar S_W)}
{\sum_{n\in W_{\ell}}(n-\bar n)^2},
\qquad
\widehat S_{i,\ell+h}^{\mathrm{OLS}}
=\bar S_{i,\ell}+h k_{i,\ell}^{\mathrm{OLS}}.
\end{equation}
Persistence基线令$\widehat S_{i,\ell+h}=\bar S_{i,\ell}$。

SG导数基线则在每个中心$n$的11圈局部窗口上拟合二阶多项式
\begin{equation}
(\widehat\alpha_0,\widehat\alpha_1,\widehat\alpha_2)
=\arg\min_{\alpha_0,\alpha_1,\alpha_2}
\sum_{j=-5}^{5}\left[S_{i,n+j}-(\alpha_0+\alpha_1j+\alpha_2j^2)\right]^2.
\end{equation}
由于$\mathrm d(\alpha_0+\alpha_1j+\alpha_2j^2)/\mathrm dj|_{j=0}=\alpha_1$，$\widehat\alpha_1$ 即局部一阶导数；取末 20 圈导数中位数 $d_{i,\ell}$，得到 SG 导数基线 $\widehat S_{i,\ell+h}^{D}=\bar S_{i,\ell}+hd_{i,\ell}$。

候选特征包括 SOH 的末端稳健水平、AUC、全局与近期斜率及一、二阶导数中位数；在此基础上逐步加入 IR、温度、充电时间的中位状态、斜率及其差值，并比较九种策略哑变量的额外信息。所有特征均在当前可见的第 1 至 $\ell$ 圈内重算：末端状态取最后 10 圈中位数，导数取最后 20 圈中位数，近期趋势统一取最近 30 圈。

增强模型先计算各训练电池相对 TS 基线的 50 维未来残差向量
\begin{equation}
\mathbf r_i=
\left(S_{i,\ell+1}-\widehat S_{i,\ell+1}^{\TS},\ldots,
S_{i,\ell+50}-\widehat S_{i,\ell+50}^{\TS}\right)^{\!\top}.
\end{equation}
将训练折的残差向量减去均值$\bar{\mathbf r}$并按行堆叠成矩阵$\mathbf R$，其协方差矩阵为
\begin{equation}
\mathbf C_r=\frac{1}{M-1}\mathbf R^{\top}\mathbf R.
\end{equation}
对$\mathbf C_r$作特征分解$\mathbf C_r\boldsymbol\phi_k=\lambda_k\boldsymbol\phi_k$，取最小的$K$使
\begin{equation}
\frac{\sum_{k=1}^{K}\lambda_k}{\sum_{k=1}^{50}\lambda_k}\ge0.95.
\end{equation}
于是残差可近似为前 $K$ 个主成分的线性组合，得分定义为
\begin{equation}
\mathbf r_i\approx\bar{\mathbf r}+\sum_{k=1}^{K}z_{ik}\boldsymbol\phi_k,
\qquad
z_{ik}=(\mathbf r_i-\bar{\mathbf r})^{\top}\boldsymbol\phi_k.
\end{equation}

用Ridge回归\cite{hoerl1970}从标准化早期特征矩阵$\mathbf X$预测得分矩阵$\mathbf Z$：
\begin{equation}
J(\mathbf B)=\|\mathbf Z-\mathbf X\mathbf B\|_F^2
+\lambda\|\mathbf B\|_F^2.
\end{equation}
对$\mathbf B$求导并令其为0：
\begin{align}
\frac{\partial J}{\partial\mathbf B}
&=-2\mathbf X^{\top}(\mathbf Z-\mathbf X\mathbf B)+2\lambda\mathbf B=0,\\
(\mathbf X^{\top}\mathbf X+\lambda\mathbf I)\mathbf B
&=\mathbf X^{\top}\mathbf Z,\\
\widehat{\mathbf B}
&=(\mathbf X^{\top}\mathbf X+\lambda\mathbf I)^{-1}\mathbf X^{\top}\mathbf Z.
\end{align}
对新电池，得分预测为 $\widehat{\mathbf z}_i=\mathbf x_i^{\top}\widehat{\mathbf B}$，残差预测为 $\widehat{\mathbf r}_i=\bar{\mathbf r}+\sum_{k=1}^{K}\widehat z_{ik}\boldsymbol\phi_k$，最终预测等于 TS 基线加残差修正。特征标准化、PCA 维数 $K$ 与正则参数 $\lambda$ 均只在每个 LOBO 训练折内确定。

\subsubsection{模型选择、消融与信息贡献}

在 $\ell=150$ 的直接基线比较中，Persistence、OLS、TS 与 SG 导数基线的 RMSE 依次为 0.002233、0.0005343、0.0005289 和 0.0006730，TS 仅比 OLS 低约 1.02\%。增强模型中 M2 的 RMSE 最低（0.000583），仍未超过 TS；M3 加入温度和充电时间后 RMSE 增至 0.001882，M4 加入策略标签后回落至 0.000598，说明这些特征虽能解释部分残差，却未形成稳定的样本外净收益。

对任一内层候选$m$，以电池为独立单位计算$e_{im}=\mathrm{MSE}_{im}$，并定义
\begin{equation}
\bar e_m=\frac1M\sum_{i=1}^{M}e_{im},
\qquad
\mathrm{SE}_m=\frac{\mathrm{sd}(e_{1m},\ldots,e_{Mm})}{\sqrt M}.
\end{equation}
若 $m^*=\arg\min_m\bar e_m$，则保留满足 $\bar e_m\le\bar e_{m^*}+\mathrm{SE}_{m^*}$ 的近优集合，并选择其中复杂度最低者。OLS 与 TS 同为单斜率趋势基线，二者同时进入近优集合时按预设的稳健基线规则优先 TS。最终 40 个外层折均选择 TS\_ONLY。

可见对第 151--200 圈预测，最可靠的信息是电池自身最近 30 圈的 SOH 退化趋势，策略标签、IR、温度与充电时间等特征均未形成稳定的增量收益。嵌套 LOBO 的最终结果为：总体 $\MAE=0.000343$、$\RMSE=0.000529$、$R^2=0.9970$，相对吻合度 99.965\%，第 200 圈绝对误差中位数 0.000395。

固定 Theil--Sen 基线在不同可见长度下的结果进一步表明，RMSE 随 $\ell$ 由 50 增至 150 圈从 0.00406 降至 0.000529，说明第 100 圈之后的近期趋势对随后 50 圈预测尤为重要；正式任务取 $\ell=150$ 由题设可见数据长度确定，不产生额外选择偏差。

\subsubsection{测试电池预测与EOL外推}

将最终趋势模型用于 9 块测试电池，得到第 151--200 圈 SOH 预测。各预测曲线在第 150 圈边界连续，并对快速退化的 16 号电池给出了明显更陡的下降趋势。

将每块测试电池的前 150 圈观测与未来 50 圈预测合并，并严格继承问题一的 $n_0=31$ 中心化单调二次结构：
\begin{equation}
S_i(n)=A_i-B_i(n-31)-C_i(n-31)^2,\qquad B_i,C_i\ge0.
\end{equation}
令$x_i=\widehat L_i^{\mathrm{ind}}-31$，由
\begin{align}
A_i-B_ix_i-C_ix_i^2&=0.8,\\
C_ix_i^2+B_ix_i-(A_i-0.8)&=0
\end{align}
得到个体交点
\begin{equation}
\widehat L_i^{\mathrm{ind}}=31+
\frac{-B_i+\sqrt{B_i^2+4C_i(A_i-0.8)}}{2C_i},\qquad C_i>0.
\end{equation}

为抑制少数电池的曲率外推漂移，对 9 块测试电池引入同策略部分池化。设 $\widehat L_{p(i)}^{\mathrm{peer}}$ 为训练折内同策略电池 EOL 的几何中心，定义
\begin{align}
\log \widehat L_i^{\mathrm{final}}
&=w\log \widehat L_i^{\mathrm{ind}}
+(1-w)\log \widehat L_{p(i)}^{\mathrm{peer}},\\
\widehat L_i^{\mathrm{final}}
&=(\widehat L_i^{\mathrm{ind}})^w
\cdot(\widehat L_{p(i)}^{\mathrm{peer}})^{1-w}.
\end{align}
权重 $w$ 在各外层训练折内选择，40/40 折均取 $w=0.75$，故正式点估计为 75\% 个体信息与 25\% 同策略训练信息的几何收缩；该步骤不回填问题一的 49 块统一基准寿命。

为传播第 151--200 圈短期预测误差，按 Bootstrap 思想\cite{efron1979}以"电池块"为单位，从 40 块训练电池的整段未来残差中有放回抽样 2000 次，并对每次交点施加相同的同策略收缩，以保留同一电池 50 圈残差的时间相关性。部分池化使 EOL 截断更新中位差由 11.77\% 降至 6.72\%、均值由 30.66\% 降至 17.99\%，主要作用是稳健化尾部漂移。

近期斜率变化辅助头在四个候选中按 one-SE 简约原则取全体均值 $\widehat{\Delta k}$，使未来斜率 LOBO RMSE 较直接延续改善 24.99\%；该头仅说明近期退化存在总体加速修正，不参与 EOL 结构门控。

\subsubsection{长期三均值辅助压力测试}

题目同时给出每块电池完整寿命区间的平均内阻、平均温度和平均充电时间。为避免长期信息泄漏，先完全冻结 SOH 预测与 EOL，再仅用 40 块训练电池前 150 圈的逐循环数据拟合状态关系。在 Policy 交互、全局锚定与 Persistence 三种关系的电池级 LOBO 比较中，三者均落入最优误差的一标准误范围，故正式关系全部采用 Persistence，避免复杂关系带来样本外波动。

将已观测的 1--150 圈与冻结关系生成的 151--$\widetilde L_i$ 圈拼接，计算完整寿命均值。与题目真实汇总量相比，IR、温度、充电时间三通道吻合度分别为 98.74\%、98.25\% 与 95.65\%，但该指标衡量的是完整寿命状态均值而非 EOL 圈数准确率。结合"预测至 200 圈/真实观察至 200 圈"同结构外推中位差 17.23\%、排序相关仅 0.654 的结构审计，本文的结论是：第 151--200 圈短期 SOH 预测获得了直接标签支持，长期 EOL 仍缺少真实标签与可靠的跨通道验证。

\subsection{问题四：充电策略优化}

\subsubsection{优化目标、时间模型与主退化指标}

令策略向量为$\boldsymbol\theta=(C_1,q,C_2)$。若恒流倍率为$C$、充入SOC比例为$\Delta s$，由"时间=容量/电流"得到$60\Delta s/C$分钟，故两阶段相加为
\begin{align}
T_1(\boldsymbol\theta)&=60\frac{q}{C_1},\\
T_2(\boldsymbol\theta)&=60\frac{0.8-q}{C_2},\\
T_{0\rightarrow80}^{\mathrm{theo}}(\boldsymbol\theta)
&=T_1+T_2
=60\left(\frac{q}{C_1}+\frac{0.8-q}{C_2}\right).
\end{align}
逐循环\texttt{chargetime}按前述口径记作$T_{i}^{\mathrm{obs}}$。已有策略$g$的观测时间与退化分别汇总为
\begin{equation}
T_g^{\mathrm{obs}}=\operatorname{median}_{i\in g}T_i^{\mathrm{obs}},
\qquad
R_g=\operatorname{median}_{i\in g}r_i^{\mathrm{stable}}.
\end{equation}
概念上的双目标为
\begin{equation}
\min_{\boldsymbol\theta}\left(
T_{0\rightarrow80}(\boldsymbol\theta),
r^{\mathrm{stable}}(\boldsymbol\theta)
\right).
\end{equation}
但Q1/Q3的80\% EOL没有真实标签，且"预测至200圈/真实观察至200圈"后的同结构外推中位差仍为17.23\%、排序相关仅0.654，故问题四不把EOL点估计作为连续优化主目标，而以直接观测的$r^{\mathrm{stable}}$表示寿命损伤。已有策略比较使用$T_g^{\mathrm{obs}}$；新参数没有实测时间，只把物理公式作为设计约束，并单列其理论--实测偏差风险。

\subsubsection{九种已有策略的经验Pareto优选}

对两个均需最小化的指标，若策略$g_1$满足
\begin{equation}
T_{g_1}^{\mathrm{obs}}\le T_{g_2}^{\mathrm{obs}},\qquad
R_{g_1}\le R_{g_2},
\end{equation}
且至少一个不等号严格成立，则称$g_1$支配$g_2$。不被任何其他策略支配的集合为
\begin{equation}
\mathcal P_{\mathrm{obs}}
=\left\{g:\nexists h\ne g, h\prec g\right\}.
\end{equation}
表中使用全部49块电池前150圈数据。为评估小样本下前沿身份的波动，另在每个策略内部有放回重采样5000次，并统计每次进入Pareto集的频率；该频率是稳定性描述，不是显著性$p$值。

精确支配下，N2（5.3C--54\%--4.0C新结构）、N5（5.6C--36\%--4.3C新结构）与N6（4.8C--80\%--4.8C新结构）构成三点前沿。N5比N2仅快约0.009 s，但观测退化速率高21.6\%，这一时间差不具工程分辨意义；N6的观测退化速率比N2低18.8\%，代价是约0.995 min更长的记录时间。因此，在强调约10分钟快充时N2是更合理的已有策略；若可接受约1 min额外时间并优先压低直接观测退化，N6是备选。

\subsubsection{六个近等时新结构策略的局部可行域}

连续优化只使用六个同时属于数据集3和新结构类别的参数位置。记其原始参数点为$\boldsymbol\theta_j$，$j=1,\ldots,6$，先定义原始参数凸包
\begin{equation}
\Omega_{\theta}=\operatorname{conv}\{\boldsymbol\theta_1,\ldots,\boldsymbol\theta_6\}.
\end{equation}
为避免"原始参数看似插值，但进入Q2回归后在$(A,D)$空间外推"，再令$\Phi(\boldsymbol\theta)=(A(\boldsymbol\theta),D(\boldsymbol\theta))$并要求
\begin{equation}
\Phi(\boldsymbol\theta)\in
\Omega_{AD}=\operatorname{conv}\{\Phi(\boldsymbol\theta_1),\ldots,\Phi(\boldsymbol\theta_6)\}.
\end{equation}
对原始参数按六点均值$\boldsymbol\mu_\theta$和标准差$\boldsymbol s_\theta$逐维标准化：
\begin{equation}
\boldsymbol z(\boldsymbol\theta)
=\operatorname{diag}(\boldsymbol s_\theta)^{-1}
(\boldsymbol\theta-\boldsymbol\mu_\theta).
\end{equation}
候选到最近真实策略的距离及允许上限定义为
\begin{align}
d(\boldsymbol\theta)&=\min_j
\|\boldsymbol z(\boldsymbol\theta)-\boldsymbol z(\boldsymbol\theta_j)\|_2,\\
d_{\max}&=\max_j\min_{k\ne j}
\|\boldsymbol z(\boldsymbol\theta_j)-\boldsymbol z(\boldsymbol\theta_k)\|_2=3.326.
\end{align}
于是局部可行域为
\begin{equation}
\boxed{\Omega_{\mathrm{local}}=\left\{
\boldsymbol\theta\in\Omega_{\theta}:
\Phi(\boldsymbol\theta)\in\Omega_{AD},\ d(\boldsymbol\theta)\le d_{\max}
\right\}}.
\end{equation}
采用可实际执行的分辨率$C_1,C_2\in0.1\mathbb Z$、$q\in0.01\mathbb Z$，并取六个实验设计中最大的理论时间为预算：
\begin{equation}
\mathcal G=\left\{\boldsymbol\theta\in\Omega_{\mathrm{local}}:
T_{0\rightarrow80}^{\mathrm{theo}}(\boldsymbol\theta)\le10.0132\ \mathrm{min}
\right\}.
\end{equation}
双重凸包、近邻距离和时间约束后共保留597个网格点。该集合只允许局部插值，不能扩展为任意倍率组合。

\subsubsection{Bootstrap分位数目标与独立稳定性审计}

问题二模型使用标准化$A^*,D^*$。在第$b$次Bootstrap中，对每个策略内部按电池有放回重采样，重新计算六个策略中位退化速率并拟合
\begin{equation}
\widehat r^{(b)}(\boldsymbol\theta)
=\widehat\beta_0^{(b)}+
\widehat\beta_A^{(b)}A^*(\boldsymbol\theta)+
\widehat\beta_D^{(b)}D^*(\boldsymbol\theta).
\end{equation}
不用均值最漂亮的点，而以较不利侧90\%分位数为选择目标：
\begin{equation}
J_{0.90}(\boldsymbol\theta)
=Q_{0.90}\left(\widehat r^{(1)}(\boldsymbol\theta),\ldots,
\widehat r^{(5000)}(\boldsymbol\theta)\right),
\qquad
\boldsymbol\theta^*=\arg\min_{\boldsymbol\theta\in\mathcal G}J_{0.90}(\boldsymbol\theta).
\end{equation}
为避免在同一批重采样上选择又评价，另生成独立5000次Bootstrap作稳定性审计。令$\boldsymbol\theta_E$为选择阶段最佳已有策略N2，定义配对差
\begin{equation}
\Delta^{(b)}=\widehat r^{(b)}(\boldsymbol\theta^*)-
\widehat r^{(b)}(\boldsymbol\theta_E).
\end{equation}
只有同时满足$Q_{0.90}(\Delta)<0$，且中位改善超过Q2策略级LOPO MAE，才把新点升级为正式推荐；这一裁决在查看验证结果前固定。

局部搜索得到$\boldsymbol\theta^*=(5.1\mathrm C,45\%,4.5\mathrm C)$。独立验证中，N2与局部候选的预测分位数对比显示：候选的边际$Q_{0.90}$较低，但配对比较显示其中位改善为$-2.43\times10^{-7}$（即候选中位预测反而略高）；$P(\Delta<0)=47.1\%$，$Q_{0.90}(\Delta)=3.24\times10^{-6}>0$。这些差异也没有超过LOPO MAE $7.91\times10^{-6}$，故不满足新策略推荐门槛。

\subsubsection{最终推荐及与典型长短寿命策略比较}

综合经验Pareto、独立Bootstrap和LOPO可分辨尺度，\paperstrong{最终推荐已有实验策略N2：5.3C--54\%--4.0C新结构}。局部候选C$^*$（5.1C--45\%--4.5C）仅作为后续实验候选，不应在缺少真实循环数据时宣称优于N2。该推荐对照下，N6的观测退化速率比N2低约16.5\%，且基准EOL更高，但观测充电时间约多0.995 min；因此在强调约10分钟快充时N2更合适，若允许约1 min额外时间并优先压低退化，N6是备选。

\subsubsection{问题三短期一致性与适用风险}

对40块完整电池，取外层OOF在第200圈的预测，并计算第151--200圈观测与预测平均退化速率。九策略的$r_g^{\mathrm{stable}}$与真实未来退化速率的Spearman相关为0.750（$p=0.0199$），与OOF预测未来退化速率的相关为0.833（$p=0.00527$）；样本仅九个策略，故这里只作方向一致性审计。N2在真实未来退化速率中排第2，N6排第1，均未与问题四的低退化判断发生明显冲突。

全新策略C$^*$没有前150圈SOH，不能直接输入问题三。合理部署流程是：先把C$^*$作为实验候选运行早期循环；获得至少150圈真实SOH后，再用问题三最近30圈Theil--Sen模型更新第151--200圈短期健康预测。最终结论受以下边界限制：连续响应面只有六个同类参数位置，局部候选仍可能受未建模交互和批次效应影响；理论时间适合约束设计，但4.8C--80\%--4.8C新结构存在1.038 min理论--实测偏差，新候选的真实时间必须实验确认；\texttt{NEWSTRUCTURE}只是数据类别标签，不能外推为所有新结构电芯的普遍规律；推荐主要降低前150圈稳定退化速率，不等同真实80\% EOL已经得到优化；N2是当前数据内的主推荐，N6是允许约1 min额外时间时的直接观测低退化备选，C$^*$只适合进一步实验筛查。

\section{模型评价与结论边界}

\subsection{优点}

本文建模方法的优势主要体现在四个方面。其一，预处理仅修复有充分局部证据的孤立尖峰，避免误删真实退化；第 1--150 圈建模与第 151--200 圈验证严格隔离，Q1 与 Q3 均以嵌套 LOBO 评价完整选型流程，防止未来信息参与模型选择。其二，Q1 将稳定起点与线性、幂律、中心化二次三类结构开放选择；Q3 的同策略池化权重在 40 个外层折中一致取 0.75，同时降低 EOL 更新差的中位数与均值；参数效应以策略为统计单位并限定在同一结构、同一数据集内，避免伪重复与批次混杂。其三，增强模型未取得超过一标准误的稳定收益时保留简单模型，Q3 的 40 折均选择 TS\_ONLY，结果简洁且可复现。其四，Q4 先基于已有实验策略给出经验 Pareto 参照，再以原始参数与 $(A,D)$ 双重凸包限制连续搜索，并用独立 Bootstrap 与 LOPO 误差尺度阻止不稳定的新策略推荐。

\subsection{局限与后续工作}

本文工作的局限在于：其一，观测 SOH 大多远高于 0.8，寿命圈数对微小斜率误差高度敏感，绝对 EOL 只能作为基准外推；九种策略跨三个数据集且部分策略仅 3 块电池，跨数据集排序可能混杂批次与结构影响。其二，数据集 3 仅有六个不同参数位置、残差自由度仅 3，主响应精确置换虽达 $p=0.0181$ 且 LOPO 方向稳定，仍不能作全参数空间的因果识别；EOL 辅助回归的 Bootstrap 区间跨 0，仅作方向参考。其三，Q3 的寿命 Bootstrap 只传播未来 50 圈经验残差，未覆盖退化函数的结构不确定性，故 Q4 以直接观测的稳定退化速率而非 EOL 点估计作主优化目标；嵌套 LOBO 仅为题设附件内部验证，$R^2=0.9970$ 只支持第 151--200 圈短期预测，不代表真实 EOL 精度。其四，完整寿命三均值的吻合度不等于 EOL 准确率，两种同结构外推的中位差仍为 17.23\%、Spearman 仅 0.654；Q4 局部候选的时间与退化均来自模型，仍须以新的电池实验验证后方可用于实际充电控制。

综上，本文以前 150 圈数据能够区分明显的快、慢退化策略，并以约 99.965\% 的 SOH 相对吻合度预测随后 50 圈；在六个近等时新结构策略位置上，较高平均倍率和更偏向中高 SOC 的倍率暴露与更快稳定退化相关。九种真实策略的时间--退化 Pareto 前沿包含 N2、N5 和 N6，约 10 分钟预算下最终推荐已有实验策略 N2（5.3C--54\%--4.0C 新结构）。局部候选未通过独立配对 Bootstrap 与 LOPO 尺度检验，仅作待实验方案。Q1/Q3 的中心化二次 EOL 经部分池化后内部稳定性提高，但仍为一致口径下的情景指标；问题四的推荐降低的是早期稳定退化风险，尚不能等同真实寿命已被精确优化。

\begin{thebibliography}{99}
\bibitem{severson2019}
Severson K A, Attia P M, Jin N, et al. Data-driven prediction of battery cycle life before capacity degradation[J]. Nature Energy, 2019, 4: 383--391. DOI: 10.1038/s41560-019-0356-8.

\bibitem{savitzky1964}
Savitzky A, Golay M J E. Smoothing and differentiation of data by simplified least squares procedures[J]. Analytical Chemistry, 1964, 36(8): 1627--1639. DOI: 10.1021/ac60214a047.

\bibitem{sen1968}
Sen P K. Estimates of the regression coefficient based on Kendall's tau[J]. Journal of the American Statistical Association, 1968, 63(324): 1379--1389. DOI: 10.1080/01621459.1968.10480934.

\bibitem{kruskal1952}
Kruskal W H, Wallis W A. Use of ranks in one-criterion variance analysis[J]. Journal of the American Statistical Association, 1952, 47(260): 583--621. DOI: 10.1080/01621459.1952.10483441.

\bibitem{dunn1964}
Dunn O J. Multiple comparisons using rank sums[J]. Technometrics, 1964, 6(3): 241--252. DOI: 10.1080/00401706.1964.10490181.

\bibitem{holm1979}
Holm S. A simple sequentially rejective multiple test procedure[J]. Scandinavian Journal of Statistics, 1979, 6(2): 65--70.

\bibitem{hoerl1970}
Hoerl A E, Kennard R W. Ridge regression: biased estimation for nonorthogonal problems[J]. Technometrics, 1970, 12(1): 55--67. DOI: 10.1080/00401706.1970.10488634.

\bibitem{efron1979}
Efron B. Bootstrap methods: another look at the jackknife[J]. The Annals of Statistics, 1979, 7(1): 1--26. DOI: 10.1214/aos/1176344552.

\bibitem{ansean2016}
Anse\'an D, Dubarry M, Devie A, et al. Fast charging technique for high power LiFePO$_4$ batteries: A mechanistic analysis of aging[J]. Journal of Power Sources, 2016, 321: 201--209. DOI: 10.1016/j.jpowsour.2016.04.140.

\bibitem{gao2017}
Gao Y, Jiang J, Zhang C, et al. Lithium-ion battery aging mechanisms and life model under different charging stresses[J]. Journal of Power Sources, 2017, 356: 103--114. DOI: 10.1016/j.jpowsour.2017.04.084.
\end{thebibliography}

\newpage
\section*{附录：复现说明}

\subsection*{A. 逐电池关键结果}

\begin{longtable}{rrlrrrr}
\caption{49块电池的策略、充电时间、SOH与基准估计寿命}\label{tab:q1battery}\\
\toprule
电池 & 数据集 & 快充策略 & $T_{1:150}$/min & $S_{150}$ & $10^4k^{\TS}$ & $L$/圈 \\
\midrule
\endfirsthead
\multicolumn{7}{c}{表\thetable\ （续）}\\
\toprule
电池 & 数据集 & 快充策略 & $T_{1:150}$/min & $S_{150}$ & $10^4k^{\TS}$ & $L$/圈 \\
\midrule
\endhead
\midrule\multicolumn{7}{r}{续下页}\\\endfoot
\bottomrule\endlastfoot
1 & 1 & 3.6C--80\%--3.6C & 13.358 & 1.00303 & -0.197 & 1070 \\
2 & 1 & 3.6C--80\%--3.6C & 13.342 & 0.99533 & -0.200 & 1709 \\
3 & 1 & 3.6C--80\%--3.6C & 13.342 & 0.99624 & -0.176 & 1389 \\
4 & 2 & 80\%--3.6C & 13.374 & 0.97758 & -1.586 & 629 \\
5 & 2 & 80\%--3.6C & 13.342 & 0.99638 & -0.272 & 1114 \\
6 & 2 & 80\%--3.6C & 13.342 & 0.99634 & -0.268 & 1548 \\
7 & 2 & 4.8C--80\%--4.8C & 10.054 & 0.99511 & -0.406 & 808 \\
8 & 2 & 4.8C--80\%--4.8C & 10.093 & 0.99021 & -0.849 & 612 \\
9 & 2 & 4.8C--80\%--4.8C & 10.021 & 0.98955 & -0.938 & 640 \\
10 & 3 & 5.0C--67\%--4.0C（新结构） & 10.042 & 0.99721 & -0.231 & 2843 \\
11 & 3 & 5.3C--54\%--4.0C（新结构） & 10.041 & 0.99682 & -0.267 & 7638 \\
12 & 3 & 5.6C--19\%--4.6C（新结构） & 10.044 & 0.99641 & -0.282 & 7144 \\
13 & 3 & 5.6C--36\%--4.3C（新结构） & 10.041 & 0.99696 & -0.255 & 7894 \\
14 & 3 & 5.6C--19\%--4.6C（新结构） & 10.046 & 0.99765 & -0.204 & 1962 \\
15 & 3 & 5.6C--36\%--4.3C（新结构） & 10.042 & 0.99616 & -0.305 & 1434 \\
16 & 3 & 3.7C--31\%--5.9C（新结构） & 10.096 & 0.97277 & -1.925 & 856 \\
17 & 3 & 4.8C--80\%--4.8C（新结构） & 11.039 & 0.99833 & -0.127 & 4605 \\
18 & 3 & 5.0C--67\%--4.0C（新结构） & 10.043 & 0.99710 & -0.278 & 1340 \\
19 & 3 & 5.3C--54\%--4.0C（新结构） & 10.042 & 0.99845 & -0.135 & 2127 \\
20 & 3 & 4.8C--80\%--4.8C（新结构） & 11.038 & 0.99771 & -0.216 & 1659 \\
21 & 3 & 5.6C--19\%--4.6C（新结构） & 10.043 & 0.99582 & -0.368 & 1234 \\
22 & 3 & 5.6C--36\%--4.3C（新结构） & 10.043 & 0.99797 & -0.176 & 4155 \\
23 & 3 & 5.6C--19\%--4.6C（新结构） & 10.042 & 0.99604 & -0.369 & 1261 \\
24 & 3 & 5.6C--36\%--4.3C（新结构） & 10.042 & 0.99767 & -0.213 & 1305 \\
25 & 3 & 4.8C--80\%--4.8C（新结构） & 11.040 & 0.99730 & -0.230 & 8754 \\
26 & 3 & 5.3C--54\%--4.0C（新结构） & 10.042 & 0.99601 & -0.222 & 1313 \\
27 & 3 & 5.6C--19\%--4.6C（新结构） & 10.043 & 0.99574 & -0.348 & 5816 \\
28 & 3 & 5.6C--36\%--4.3C（新结构） & 10.043 & 0.99572 & -0.357 & 5648 \\
29 & 3 & 5.0C--67\%--4.0C（新结构） & 10.043 & 0.99731 & -0.250 & 1504 \\
30 & 3 & 3.7C--31\%--5.9C（新结构） & 10.091 & 0.97922 & -1.549 & 1307 \\
31 & 3 & 5.0C--67\%--4.0C（新结构） & 10.042 & 0.99760 & -0.212 & 1122 \\
32 & 3 & 5.3C--54\%--4.0C（新结构） & 10.043 & 0.99754 & -0.205 & 1133 \\
33 & 3 & 5.6C--19\%--4.6C（新结构） & 10.042 & 0.99737 & -0.230 & 1276 \\
34 & 3 & 5.6C--36\%--4.3C（新结构） & 10.045 & 0.99657 & -0.292 & 1096 \\
35 & 3 & 3.7C--31\%--5.9C（新结构） & 10.112 & 0.97188 & -2.167 & 636 \\
36 & 3 & 5.3C--54\%--4.0C（新结构） & 10.044 & 0.99753 & -0.211 & 1420 \\
37 & 3 & 5.0C--67\%--4.0C（新结构） & 10.043 & 0.99657 & -0.298 & 1579 \\
38 & 3 & 5.3C--54\%--4.0C（新结构） & 10.044 & 0.99579 & -0.357 & 1661 \\
39 & 3 & 5.6C--19\%--4.6C（新结构） & 10.044 & 0.99643 & -0.290 & 2319 \\
40 & 3 & 5.6C--36\%--4.3C（新结构） & 10.043 & 0.99540 & -0.374 & 1752 \\
41 & 3 & 4.8C--80\%--4.8C（新结构） & 11.038 & 0.95167 & -0.314 & 1295 \\
42 & 3 & 5.0C--67\%--4.0C（新结构） & 10.044 & 0.99755 & -0.200 & 1938 \\
43 & 3 & 5.3C--54\%--4.0C（新结构） & 10.044 & 0.99640 & -0.312 & 1526 \\
44 & 3 & 5.6C--19\%--4.6C（新结构） & 10.043 & 0.99654 & -0.323 & 1358 \\
45 & 3 & 5.6C--36\%--4.3C（新结构） & 10.043 & 0.99826 & -0.147 & 1302 \\
46 & 3 & 4.8C--80\%--4.8C（新结构） & 11.038 & 0.99836 & -0.146 & 1514 \\
47 & 3 & 5.0C--67\%--4.0C（新结构） & 10.044 & 0.99840 & -0.146 & 1321 \\
48 & 3 & 5.3C--54\%--4.0C（新结构） & 10.043 & 0.99829 & -0.162 & 1830 \\
49 & 3 & 4.8C--80\%--4.8C（新结构） & 11.038 & 0.99833 & -0.143 & 1891 \\
\end{longtable}


\subsection*{B. 文件结构}

\begin{itemize}
  \item \texttt{问题1（latest）(2).md}、\texttt{问题2（latest）(2).md}、\texttt{问题3（latest）(2).md}与\texttt{问题4（latest）.md}：四问的当前建模规范；仅定义口径、候选模型和裁决规则，不保存本轮实证数值；
  \item \texttt{scripts/analysis\_pipeline.py}：数据校验、异常处理、Q1--Q3基础模型、交叉验证、置换检验与Bootstrap；
  \item \texttt{scripts/eol\_v4\_final\_arbitration.py}：15候选联合裁决、Q3部分池化和未来斜率四基线比较；
  \item \texttt{scripts/formal\_v4\_pipeline.py}：从原始CSV串联Q1--Q4并执行寿命传播审计；
  \item \texttt{scripts/q4\_optimization.py}：经验Pareto、双重凸包局部网格、分离Bootstrap选择/验证和Q3一致性审计；
  \item \texttt{tests/test\_q4\_optimization.py}：时间公式、Pareto支配、局部域与Bootstrap分位数的单元测试；
  \item \texttt{scripts/make\_figures.R}：全部论文图的R代码；
  \item \texttt{results/正式重跑\_20260816\_v4/summary.json}：四问正式主管线的核心结论和数值汇总；
  \item 同目录的\texttt{问题1/问题2/问题3/问题4}：各问中间结果、交叉验证与最终预测；
  \item \texttt{results/正式重跑\_20260816\_v4/Q1-Q4寿命结果全链路一致性审计.csv}：Q1寿命向Q2、Q4及Q3部分池化的逐项传播检查；
  \item \texttt{figures\_generated\_v4/}：论文图文件。
\end{itemize}

\subsection*{C. 关键计算顺序}

\begin{enumerate}[label=\arabic*)]
  \item 读取两个题目CSV并检查电池--循环唯一性；
  \item 按电池局部Hampel/MAD标记异常，仅修复满足局部孤立性条件的容量、内阻和温度尖峰；
  \item 重算SOH并做SG(11,2)平滑；
  \item 联合比较五个稳定起点与三类低参数寿命结构，以151--200圈误差形成one-SE近优集合，再按无效EOL、截断更新、边界、复杂度和起点裁决，并用外层LOBO评价完整选型流程；
  \item 汇总九种策略，完成KW置换检验、Dunn--Holm比较和数据集3策略级参数回归；
  \item 在40块训练电池上完成外层LOBO、内层候选选择、一标准误简约决策、开发阶段特征消融和早期长度敏感性，再冻结模型预测9块测试电池；
  \item 用训练电池未来残差按电池块重采样，计算测试电池个体EOL，再以75\%个体、25\%同策略训练中心作对数尺度部分池化；
  \item 只用训练电池前150圈拟合三种状态关系并作电池级LOBO；在SOH/EOL冻结后才与三个完整寿命均值比较，且不据此反向校准寿命；
  \item 以九策略观测时间和稳定退化速率构造经验Pareto前沿；
  \item 在六个同类参数位置的原始参数凸包与$(A,D)$特征凸包交集中生成可实施网格，并施加近邻距离和10.0132 min理论时间预算；
  \item 用5000次Bootstrap选择$Q_{0.90}$最小候选，再用独立5000次Bootstrap、配对差和Q2 LOPO误差尺度决定是否升级为正式推荐；
  \item 运行绘图脚本生成全部图形，再由XeLaTeX编译论文。
\end{enumerate}

\end{document}
